% Lösungen Einheit 5 – Nr. 58–69
\einheitenkopf[le5][5 Potenzfunktionen]{Einheit 5 von 5 · Potenzfunktionen mit natürlichem Exponenten}

\erg{58}{a) 0; 1; 8; 27 \quad b) 0; 1; 16; 81 \quad c) 64; 125; 1000 \quad d) 256; 625; 10\,000 \quad e) $-27$; $-8$; $-1$ \quad f) 81; 16; 1}
\erg{59}{a) $-4$; $-0{,}5$; 0; 0,5; 4 \quad b) $-16$; $-1$; 0; $-1$; $-16$ \quad c) $-3{,}375$; $-0{,}125$; 0,125; 3,375}
\erg{60}{a) kubische Parabel durch $(-2\,|\,{-8})$, $(-1\,|\,{-1})$, $(0\,|\,0)$, $(1\,|\,1)$, $(2\,|\,8)$ \quad b) Parabel vierter Ordnung durch $(-1{,}5\,|\,5{,}06)$, $(-1\,|\,1)$, $(0\,|\,0)$, $(1\,|\,1)$, $(1{,}5\,|\,5{,}06)$, flach um den Ursprung \quad c) durch $(-2\,|\,{-4})$, $(-1\,|\,{-0{,}5})$, $(0\,|\,0)$, $(1\,|\,0{,}5)$, $(2\,|\,4)$}
\erg{61}{a) 8 \quad b) 64 \quad c) 81 \quad d) 1 \quad e) $-8$ \quad f) 16 \quad g) $-32$ \quad h) 3,375 \quad i) 0,0625}
\erg{62}{A ja \quad B nein ($2^4 = 16$) \quad C ja \quad D nein ($(-2)^4 = 16$) \quad E ja \quad F nein ($-(-1)^4 = -1$)}
\erg{63}{a) 2 \quad b) 10 \quad c) 1 \quad d) 5 \quad e) $x = 2$ oder $x = -2$ \quad f) $-3$ \quad g) keine Lösung \quad h) $x \approx 3{,}68$}
\erg{64}{a) $y$-Achse \quad b) Ursprung \quad c) $y$-Achse \quad d) Ursprung \quad e) $y = 3x^4$ \quad f) $y = -x^4$}
\erg{65}{a) B \quad b) A \quad c) D \quad d) C}
\erg{66}{a) Potenz \quad b) Exponential \quad c) Potenz \quad d) Exponential \quad e) Potenz \quad f) Exponential}
\erg{67}{a) Tim hat $-x^4$ wie $(-x)^4$ gerechnet; das Minus steht vor der Potenz: $f(2) = -2^4 = -16$ \quad b) Sara hat $x^3$ als $3 \cdot x$ gerechnet: $g(5) = 5 \cdot 5 \cdot 5 = 125$}
\erg{68}{a) Eine Zahl hoch 4 ist ein Produkt aus vier gleichen Faktoren; bei negativem $x$ ergeben je zwei Minuszeichen Plus, also ist $x^4 \ge 0$. \quad b) Für negatives $x$ ist $x^3$ ein Produkt aus drei negativen Faktoren und damit negativ: $x < 0$ und $y < 0$ ist der dritte Quadrant.}
\erg{69}{a) 125\,cm³ \quad b) 10\,cm \quad c) $k \approx 7{,}9$\,cm (7,937) \quad d) Es wird 8-mal so groß: $(2k)^3 = 8k^3$}
